\documentclass[11pt]{article}
\usepackage[utf8]{inputenc}
\usepackage{graphicx}
\usepackage{amsmath}
\usepackage{amsfonts}
\usepackage{amssymb}
\usepackage{listings}
\usepackage{xcolor}
\usepackage{hyperref}
\usepackage{geometry}
\geometry{margin=1in}

% Code listing style
\lstset{
    basicstyle=\ttfamily\small,
    breaklines=true,
    frame=single,
    numbers=left,
    numberstyle=\tiny\color{gray},
    keywordstyle=\color{blue},
    commentstyle=\color{green!60!black},
    stringstyle=\color{red}
}

\title{Automated Data Extraction and Validation Pipeline for Global Energy Storage Database}
\author{Your Name/Organization}
\date{\today}

\begin{document}

\maketitle

\begin{abstract}
This technical report describes an automated pipeline for extracting, processing, and validating energy storage project data from web sources for the Global Energy Storage Database (GESDB). The system combines web scraping, large language model-based information extraction, comprehensive validation, and a web-based interface for data management. The pipeline automates the collection of project information from online sources, extracts structured data using a modular tool-based architecture, validates extracted data against schema rules, and provides a user interface for review and integration. This approach significantly reduces manual data entry effort while maintaining data quality through systematic validation processes.
\end{abstract}

\section{Introduction}
\label{sec:introduction}

The Global Energy Storage Database (GESDB) requires comprehensive and accurate information about energy storage projects worldwide, including technical specifications, location data, project status, ownership, and operational characteristics. Manually collecting and maintaining this information from diverse web sources is time-consuming and error-prone. This technical report presents an automated pipeline that addresses these challenges by systematically extracting structured data from web articles and validating it against established database schema requirements.

The pipeline consists of five main stages: (1) web scraping to collect article content from energy storage project websites, (2) information extraction using large language models to convert unstructured text into structured data fields, (3) optional data autofill to enhance completeness by searching the web for missing values, (4) comprehensive validation against schema rules with detailed error reporting, and (5) web-based streamlit interface for data review, editing, and integration. The system is designed with a modular architecture that allows for easy extension and customization of extraction capabilities.

The primary contributions of this work include a tool-based extraction system that enables flexible data extraction across multiple domains, a comprehensive validation framework with detailed error codes for quality assurance, and an integrated web interface that streamlines the data management workflow. The system has been applied to extract data from power-technology.com and can be extended to additional sources as needed.

\section{System Architecture and Pipeline Overview}
\label{sec:architecture}

The GESDB Pipeline follows a sequential five-stage architecture designed to transform unstructured web content into validated, structured database records. The system can be executed through an interactive web interface (via \texttt{pipeline.py}).

\subsection{Pipeline Stages}

The pipeline consists of five distinct stages that process data sequentially:

\textbf{Stage 1: Article Collection.} The system begins by scraping web sources for energy storage project articles. The scraper identifies article URLs using configurable CSS selectors, extracts article content (title, body text, and source URL) using the newspaper library, and saves the raw content to JSON files. The system maintains a record of processed URLs to avoid duplicate processing on subsequent runs.

\textbf{Stage 2: Information Extraction.} Raw article content is processed through a large language model (LLM)-based extraction system. The ArticleProcessor orchestrates multiple specialized extraction tools, each designed to extract specific categories of information such as project details, location data, dates, applications, grid interconnection information, ownership, and subsystem specifications. The extracted data is then transformed into a standardized format matching the GESDB schema.

\textbf{Stage 3: Data Autofill.} As an optional post-processing step, the system can enhance data completeness by identifying missing or null values in processed records and attempting to fill them through web search and extraction. For each missing field, the system constructs contextual search queries using project name, location, and field-specific information, then uses Google Custom Search Engine and language models to extract values from web sources. This stage generates suggestions with confidence scores that users can review and selectively merge, helping to improve data completeness while maintaining quality control.

\textbf{Stage 4: Validation.} Extracted data undergoes comprehensive validation against schema rules defined in a CSV configuration file. The validation system checks data types, value ranges, required field presence, and field-specific constraints (e.g., valid status values, application categories, technology classifications). Validation results are reported with detailed error codes, enabling precise identification of data quality issues.

\textbf{Stage 5: Data Integration.} The final stage performs data matching and merging operations. Using fuzzy matching techniques, the system identifies similar records across different data sources and merges complementary information to create comprehensive records. This stage also provides functionality for exporting validated records to external systems, such as GitHub repositories.

\subsection{Data Flow}

Data flows through the system as follows: web articles are collected and stored as JSON files in the \texttt{data/input/} directory. These articles are then processed by the extraction system, which outputs structured records to \texttt{data/output/processed\_data.json}. The validation system reads this processed data, performs validation checks, and generates validation status reports in \texttt{validation/output/}. The web interface provides access to all stages of this pipeline, allowing users to execute individual stages, review results, edit records, and manage the validation workflow.

\subsection{Component Organization}

The system is organized into modular components: the \texttt{article\_scraper} module handles web scraping, the \texttt{article\_processor} module orchestrates extraction, the \texttt{extractor} package contains the tool-based extraction system, the \texttt{postprocessor} module transforms extracted data into the target schema, the \texttt{validation} package performs schema validation, and the \texttt{pipeline.py} module provides the web interface. This modular design enables independent development, testing, and extension of individual components.

\section{Backend Components}
\label{sec:backend}

The backend of the GESDB Pipeline consists of five main components that work together to transform web articles into validated database records. These components are executed sequentially through the pipeline function in \texttt{main.py}, with each component handling a specific aspect of the data processing workflow.

\subsection{Article Scraper}

The Article Scraper component (\texttt{article\_scraper.py}) is responsible for collecting raw article content from web sources. It operates in three phases: URL discovery, content extraction, and data persistence.

The scraper begins by reading a list of previously processed URLs to avoid duplicate processing. It then fetches the target webpage and uses BeautifulSoup to parse the HTML, identifying article links through configurable CSS selectors. URLs are filtered based on an optional prefix pattern and checked against the existing URL list to ensure only new articles are processed.

For each new article URL, the scraper uses the newspaper library to extract the article's title, body text, and source URL. The newspaper library handles the complexities of parsing various website structures and extracting clean text content. Articles that fail to extract critical information (such as missing titles or empty body text) are skipped with appropriate logging.

The extracted articles are saved as JSON files in the \texttt{data/input/} directory, with each article represented as a dictionary containing title, text, and source fields. The scraper maintains separate files for tracking processed URLs, enabling incremental updates on subsequent runs.

\subsection{Article Processor}

The Article Processor (\texttt{article\_processor.py}) orchestrates the information extraction process. It serves as the central coordinator that loads articles, manages processing state, and coordinates between the extraction tools and post-processing components.

The processor maintains state information including a list of processed URLs to avoid re-processing articles, a processing limit for testing or incremental processing, and error handling parameters. When processing articles, it loads the JSON file containing scraped articles, iterates through each article, and for each unprocessed article, formats the content into a prompt string containing the article's title, source URL, and text.

The formatted prompt is passed to the Extractor system, which applies all registered extraction tools to extract structured information. The extracted data is then passed to the PostProcessor for transformation into the target schema format. Upon successful processing, the article URL is marked as processed and the structured record is appended to the output JSON file.

The processor includes error handling mechanisms that allow it to continue processing subsequent articles even if individual articles fail, with configurable limits on the number of errors before stopping. This design ensures robustness when processing large batches of articles with varying quality and structure.

\subsection{Extractor System}

The Extractor system (\texttt{extractor/extractor.py}) provides a modular, tool-based architecture for information extraction using large language models. The system uses LangChain's ChatOpenAI interface to interact with GPT-5, configured with temperature 0.0 for deterministic outputs and the Responses API for structured tool calling.

The Extractor class maintains a registry of extraction tools, where each tool is designed to extract a specific category of information. The system currently includes nine specialized tools: project information (name, power rating, capacity, status), location information (country, city, state, coordinates), date information (announced, constructed, commissioned dates), project applications (energy services categories), grid and utility information (interconnection level, ISO/RTO), project participants (developers, installers, technology providers), ownership and funding information, contact information, and subsystem specifications (storage devices, power conversion systems, balance of system components).

Each extraction tool is implemented as a LangChain tool decorated function with a Pydantic schema defining the expected output structure. When the Extractor processes an article, it iteratively calls each registered tool, binding the tool to the language model and forcing the model to use that specific tool. The model's response is then routed through ExtractorUtils (\texttt{extractor/utils/extractor\_utils.py}), which maps each tool name to a corresponding processing function that extracts and structures the tool's output.

This modular design enables easy extension: new extraction tools can be added by creating a new tool file with a Pydantic schema, implementing the extraction function, registering it with the Extractor, and providing a corresponding utility function for processing the tool's output. The tool-based approach also allows for independent validation of each extraction category and enables selective tool usage if needed in future enhancements.

\subsection{PostProcessor}

The PostProcessor component (\texttt{postprocessor/postprocessor.py}) transforms the extracted data from the Extractor's tool-based format into the standardized GESDB schema format. It performs several key functions: schema mapping, data formatting, ID assignment, and missing value handling.

The PostProcessor receives a list of extracted data dictionaries (one per article) and processes each element to create a complete database record. It maps fields from the tool-specific output structures to the target schema field names, handling variations in naming conventions and data organization. For example, it transforms the applications data from the tool's internal structure into the categorized format required by the schema, with separate lists for bulk energy services, ancillary services, transmission services, and other categories.

The component assigns unique sequential IDs to each processed record, with IDs starting from 1 and incrementing based on the number of existing records in the output file. This ensures unique identification even when processing is performed incrementally. The PostProcessor also sets a "Data Source" field to identify records processed by the pipeline.

For subsystem data, the PostProcessor formats the extracted specifications into the nested structure required by the schema, organizing information into Storage Device, Power Conversion System, and Balance of System categories. If no subsystem data is extracted, it creates a default empty subsystem structure to maintain schema compliance.

The PostProcessor also handles missing values through the PostProcessingUtils module, which applies default values and data type conversions as needed. The final output is a list of complete records ready for validation, with all required fields present (though some may contain null values if information was not available in the source article).

\subsection{Data Autofill System}

The Data Autofill System (\texttt{data\_insert.py}) is an optional post-processing component that attempts to fill missing or null values in processed records by searching the web and extracting relevant information. This system operates after the initial extraction and post-processing stages, providing a mechanism to enhance data completeness when source articles lack certain details.

The system processes each record by recursively walking through all fields in the data structure, identifying missing values (null, empty strings, empty lists or dictionaries, or NaN values for numeric fields). For each missing field, the system constructs a targeted web search query using contextual information from the record: the project name, location information (city, state, country), the field name, and field-specific descriptions defined in a configuration dictionary. This contextual approach improves search relevance by incorporating domain-specific knowledge about energy storage projects.

The search queries are executed using the Google Custom Search Engine (CSE) API, which returns relevant web page snippets. The system retrieves multiple search results (configurable, typically 5 per field) to gather diverse information sources. These snippets are then processed by a language model (GPT-4o-mini) that extracts the specific field value from the search results. The model is constrained to use only the provided snippets as evidence, returning "Unknown" if insufficient information is available, which helps maintain data quality by avoiding speculative values.

For each missing field, the system generates a suggestion containing the extracted value, a confidence score (0-1), a brief rationale explaining the extraction, and references to the source URLs. These suggestions are collected into a separate output file that does not modify the original processed data, allowing users to review and selectively merge suggestions. This design ensures that the autofill process does not introduce errors into the validated data without human oversight.

The system includes field-specific rules that guide query construction and value extraction. For example, geographic coordinates (latitude/longitude) are searched with specific formatting requirements, date fields include temporal context, and nested fields such as subsystem specifications use wildcard matching to handle array indices. The system can process nested data structures, applying the autofill logic to fields at any depth within the record hierarchy.

A significant consideration when using this system is the cost associated with third-party API calls. Each missing field requires both a Google CSE API call and an OpenAI API call, and with records containing many missing fields across multiple records, these costs can accumulate rapidly. The system provides configuration options to limit the number of search results per query and includes rate limiting capabilities (sleep intervals between queries) to manage API usage. Users should carefully consider the trade-off between data completeness and processing costs when deciding which records and fields to process through the autofill system.

\section{Validation System}
\label{sec:validation}

The validation system ensures that extracted data conforms to the GESDB schema requirements and maintains data quality standards. The validation process is executed by the \texttt{run\_validation()} function in \texttt{validation/validation\_script.py}, which performs comprehensive checks on all records in the processed data file and generates detailed validation reports.

\subsection{Validation Approach}

The validation system operates by iterating through each record in the processed data and applying multiple validation checks to each field. Validation rules are defined in a CSV configuration file (\texttt{validation/data/gesdb\_data\_rules.csv}) that specifies for each field: whether it is required or optional, the expected data type, and valid value ranges. The system uses a flexible rule lookup mechanism that handles variations in field naming (case-insensitive, spacing variations around special characters) and provides default rules for fields not explicitly defined in the configuration.

For each field, the system performs three types of validation checks: data type validation, data range validation, and field-specific validation. Data type validation ensures that values match the expected type (text, integer, float, date). Data range validation checks that numeric values fall within specified bounds and that dates are in the correct format (MM-DD-YYYY). Field-specific validation applies domain-specific constraints, such as checking that status values are from an approved list, that application categories are valid, or that technology classifications are consistent.

\subsection{Error Code System}

The validation system uses a comprehensive error code taxonomy to categorize and report validation issues. Error codes are organized into several categories, each identified by a numeric range:

\textbf{Data Type Errors (100-series):} These codes indicate issues with data types, including missing required data (101), missing optional data (102), and type mismatches for text (103), integers (104), floats (105), numeric values (106), dates (107), and generic type mismatches (108).

\textbf{Data Range Errors (200-series):} These codes report violations of value ranges, including missing data preventing range validation (201-202), lower bound violations (203-204), upper bound violations (205), and date range violations (206).

\textbf{Field-Specific Errors:} Additional error code categories handle domain-specific validation: status field errors (300-series) for invalid project status values, application field errors (700-series) for invalid application categories, URL field errors (1000-series) for malformed URLs, grid interconnection level errors (2500-series), technology classification errors (6400-series for broad categories, 6500-series for mid-types), capacity-to-power ratio warnings (20000-series), and date relationship errors (30000-series) for logical inconsistencies in project timelines.

\subsection{Validation Workflow}

The validation process begins by loading the processed data JSON file and the validation rules CSV file. For each record, the system iterates through all fields, applying validation checks based on the field's rule configuration. For complex nested structures such as Applications and Subsystems, the system recursively validates sub-fields and sub-sub-fields, ensuring comprehensive coverage of the data structure.

The validation results are aggregated for each field, with the system determining an overall validation status (Validated or Unvalidated) based on whether any validation check failed. Error codes are collected and formatted into comma-separated strings for reporting. The system also performs cross-field validation checks, such as verifying that the ratio of storage capacity to rated power is consistent with the discharge duration, and that project dates follow logical sequences (announced before constructed, constructed before commissioned, etc.).

\subsection{Validation Output}

The validation system generates multiple output files to support different use cases. The primary output is \texttt{validation/output/validation\_status.csv}, a tabular format where each row represents a record and each column represents a field, with cells containing validation status details including the validation flag, description, error codes, and timestamp. A JSON version (\texttt{validation/output/validation\_status.json}) provides the same information in a structured format suitable for programmatic access.

Additionally, the system generates \texttt{validation/output/warning\_messages.csv}, which contains warnings for cross-field validation issues such as capacity-to-power ratio anomalies and date sequence problems. These warnings do not necessarily indicate errors but flag potential data quality issues that may require manual review. The validation output is designed to be consumed by the web interface, which displays validation status alongside records and enables users to identify and correct data quality issues efficiently.

\section{Frontend Interface}
\label{sec:frontend}

The frontend interface is implemented as a Streamlit web application (\texttt{pipeline.py}) that provides an interactive environment for executing the pipeline, reviewing extracted data, managing validation workflows, and exporting validated records. The interface is organized into three main tabs, each serving a distinct purpose in the data management workflow.

\subsection{Pipeline Execution Tab}

The "Run Pipeline" tab provides a user interface for executing each stage of the data processing pipeline. The tab is organized into four expandable sections corresponding to the four pipeline stages, allowing users to execute stages individually or in sequence.

The article collection section allows users to configure scraping parameters either by using default settings or by customizing the base URL, CSS link selector, file paths, and URL prefix. When executed, the scraper runs in the background and provides status updates through the interface. The article processing section enables users to specify input and output file paths, set processing limits for testing purposes, and provide the API key for the language model. The data autofill section enables users to execute the optional autofill process to enhance data completeness by searching for missing values. The validation section provides a simple button to trigger the validation process, which reads the processed data and generates validation reports.

Each section provides appropriate warnings about processing time and displays success or error messages upon completion. This design enables users to monitor the pipeline execution, handle errors interactively, and perform incremental processing as needed.

\subsection{View Records Tab}

The "View Records" tab provides a comprehensive interface for reviewing, editing, and validating individual records. Users can select a record from a dropdown menu, and the interface displays the record's data organized into logical categories: Project Details, Location Information, Date Information, Contact Information, Grid \& Utility, Ownership and Financials, and Subsystems.

Each field within a category is displayed with an appropriate editor widget based on the data type: text inputs for strings, number inputs for numeric values, checkboxes for booleans, select boxes for fields with predefined options (such as Grid Interconnection Level and Interconnection Type), and data editors for complex nested structures like subsystems. Users can edit any field value directly in the interface, and changes are tracked in memory until explicitly saved.

The interface includes validation status indicators for each category, allowing users to mark categories as validated once they have reviewed and confirmed the data. The overall record validation status is displayed prominently, with success messages when all categories are validated. The tab also displays validation errors from the validation system, showing field-specific error codes and descriptions in an expandable section.

Users can save their edits to the processed data file, revert unsaved changes, or download the current record as a JSON file. The interface uses Streamlit's caching mechanism to efficiently load and refresh data, with a manual refresh button available to reload data from disk when needed.

\subsection{Validated Records Tab}

The "Validated Records" tab manages the workflow for exporting validated records to external systems. The tab displays a list of all records that have been marked as fully validated in the View Records tab, showing the total count and allowing users to review the JSON structure of validated records.

The primary functionality of this tab is the GitHub integration feature, which enables users to push validated records to a GitHub repository. Users can configure the target repository (owner, repository name, branch, and target directory), provide a GitHub personal access token, and optionally save the token to Streamlit's secrets file for future use. Users can select specific record IDs to export or choose to export all validated records.

When executed, the system creates individual JSON files for each record (named \texttt{record\_\{ID\}.json}) and pushes them to the specified GitHub repository using the GitHub Contents API. The interface displays the results of each push operation, indicating success or failure for each record. This integration enables seamless transfer of validated data to version-controlled repositories, supporting collaborative data management and integration with downstream systems.

\subsection{User Workflow}

The typical user workflow begins in the Pipeline Execution tab, where users run the article collection and processing stages to generate initial data. Users then navigate to the View Records tab to review the extracted data, correct any errors identified through validation, and mark categories as validated. Once records are fully validated, users move to the Validated Records tab to export the data to GitHub or other external systems.

The interface is designed to support both one-time batch processing and incremental updates. Users can process new articles, review and validate the new records, and export them without affecting previously processed data. The validation status tracking ensures that users can resume validation work across multiple sessions, with the interface maintaining state about which records and categories have been validated.

% Figure placeholders for UI screenshots
\begin{figure}[h]
    \centering
    \includegraphics[width=0.9\textwidth]{figures/pipeline_tab.png}
    \caption{Pipeline execution interface showing the four-step process}
    \label{fig:pipeline_tab}
\end{figure}

\begin{figure}[h]
    \centering
    \includegraphics[width=0.9\textwidth]{figures/view_records_tab.png}
    \caption{Record viewing and editing interface with categorized data display}
    \label{fig:view_records_tab}
\end{figure}

\begin{figure}[h]
    \centering
    \includegraphics[width=0.9\textwidth]{figures/validated_records_tab.png}
    \caption{Validated records interface with GitHub integration}
    \label{fig:validated_records_tab}
\end{figure}

\section{Integration and Workflow}
\label{sec:integration}

The frontend and backend components are integrated through direct Python function calls, with the Streamlit interface serving as a wrapper that provides user interaction and state management around the core processing functions. This design ensures that the same backend logic is used whether the system is accessed through the command-line interface (\texttt{main.py}) or the web interface (\texttt{pipeline.py}).

\subsection{Frontend-Backend Integration}

The Streamlit application imports and directly calls the backend modules: \texttt{collect\_articles} from \texttt{article\_scraper}, \texttt{ArticleProcessor} from \texttt{article\_processor}, \texttt{run\_validation} from \texttt{validation.validation\_script}, and \texttt{process\_data} from \texttt{data\_insert}. When a user clicks a button to execute a pipeline stage, the corresponding backend function is invoked with the parameters specified in the interface.

The interface uses Streamlit's caching mechanism (\texttt{@st.cache\_data}) to efficiently load the processed data file, avoiding redundant file I/O operations when the data hasn't changed. The validation status CSV file is loaded on demand when users view records, with error handling for cases where validation hasn't been run yet.

Session state management enables the interface to track validation status across page interactions. When users mark categories as validated in the View Records tab, this state is stored in Streamlit's session state and persists for the duration of the user's session. The validated records tab reads this session state to determine which records have been fully validated and are ready for export.

\subsection{Data Flow Through the System}

Data flows through the system in a unidirectional pipeline: web articles are collected and stored in \texttt{data/input/articles\_*.json}, these articles are processed by the ArticleProcessor which outputs to \texttt{data/output/processed\_data.json}, the validation system reads this processed data and generates reports in \texttt{validation/output/}, and the web interface reads both the processed data and validation reports to display them together.

The web interface can modify the processed data file directly when users save edits in the View Records tab. This creates a feedback loop where validation errors can be corrected through the interface, and the corrected data can be re-validated. The system maintains data consistency by writing changes immediately to disk and clearing caches to ensure subsequent operations see the updated data.

\subsection{Error Handling and User Feedback}

The integration includes comprehensive error handling at each stage. When backend functions raise exceptions, the Streamlit interface catches these exceptions and displays user-friendly error messages. This allows users to understand what went wrong (e.g., file not found, API errors, validation failures) and take corrective action.

The interface provides real-time feedback during long-running operations. While the backend functions execute, the interface can display progress indicators and status messages. For operations that may take significant time (such as article collection or processing), the interface displays warnings and allows users to monitor the process.

\subsection{Workflow Patterns}

The system supports several common workflow patterns. In a batch processing workflow, users execute all pipeline stages sequentially through the Pipeline Execution tab, then review and validate records in bulk. In an incremental workflow, users process new articles as they become available, validate them individually, and export validated records incrementally.

The validation workflow integrates seamlessly with the editing workflow: users can run validation to identify errors, navigate to the View Records tab to see validation errors alongside the data, make corrections directly in the interface, mark categories as validated, and export the corrected records. This integrated approach reduces the need to switch between different tools and maintains context throughout the data quality improvement process.

\section{Contributions}
\label{sec:contributions}

This work presents several key contributions to the domain of automated data extraction and validation for energy storage databases:

\textbf{Modular Tool-Based Extraction Architecture:} The system introduces a flexible, extensible architecture for information extraction using large language models. By organizing extraction capabilities into independent, registerable tools, the system enables easy addition of new extraction categories without modifying core extraction logic. Each tool uses Pydantic schemas for structured output validation, ensuring type safety and consistency. This modular design makes the system adaptable to different data domains and extraction requirements.

\textbf{Comprehensive Validation Framework:} The validation system provides a multi-level validation approach with a detailed error code taxonomy. The system validates data types, value ranges, required field presence, and domain-specific constraints, generating detailed reports that enable precise identification and correction of data quality issues. The validation rules are externalized in a CSV configuration file, allowing schema changes without code modifications. The system also performs cross-field validation to identify logical inconsistencies in the data.

\textbf{Integrated Web-Based Data Management Interface:} The Streamlit interface provides a unified environment for executing the pipeline, reviewing extracted data, correcting errors, tracking validation status, and exporting validated records. The interface integrates validation error display with data editing capabilities, enabling users to identify and correct issues in a single workflow. The category-based organization of record data improves usability by grouping related fields together, and the validation status tracking supports incremental validation workflows.

\textbf{End-to-End Automation:} The system automates the complete workflow from web scraping to validated data export, significantly reducing manual effort in data collection and entry. The pipeline handles duplicate detection, error recovery, and incremental processing, making it suitable for both initial data population and ongoing database maintenance. The automation extends to data export through GitHub integration, enabling seamless transfer of validated data to version-controlled repositories.

\textbf{Practical Application to Energy Storage Domain:} The system has been specifically designed and applied to extract energy storage project data, with extraction tools tailored to the GESDB schema requirements. The validation system includes domain-specific checks for energy storage technologies, grid interconnection levels, and project applications, demonstrating the system's applicability to real-world database management challenges.

\section{Conclusion}
\label{sec:conclusion}

This technical report has presented the GESDB Pipeline, an automated system for extracting, processing, and validating energy storage project data from web sources. The system combines web scraping, large language model-based information extraction, comprehensive validation, and an integrated web interface to streamline the data management workflow for the Global Energy Storage Database.

The pipeline's five-stage architecture (collection, extraction, autofill, validation, and integration) provides a systematic approach to transforming unstructured web content into validated, structured database records. The modular design of the extraction system enables easy extension to new data categories and sources, while the comprehensive validation framework ensures data quality through multi-level checks and detailed error reporting.

The integrated web interface bridges the gap between automated processing and human oversight, enabling users to review extracted data, correct errors identified through validation, track validation progress, and export validated records. This combination of automation and human-in-the-loop quality control addresses the challenge of maintaining accurate, comprehensive databases in domains where information is distributed across multiple web sources.

The system demonstrates the practical application of modern natural language processing and web technologies to domain-specific database management challenges. Future work could explore expanding the system to additional data sources, enhancing the extraction tools with domain-specific knowledge, and integrating with additional external systems beyond GitHub. The modular architecture provides a foundation for these extensions while maintaining the system's core functionality and usability.

% Bibliography placeholder
\begin{thebibliography}{9}
% TODO: Add references as needed
\end{thebibliography}

\end{document}

